\documentclass[12pt]{article}
\usepackage{amsmath,amssymb,amsthm}
\usepackage[margin=1in]{geometry}

\title{Problem 260: Stone Game}
\author{Project Euler Solution}
\date{}

\begin{document}
\maketitle

\section{Problem Statement}

Three piles of stones. A player may remove any positive number of stones from one pile, or an equal positive number from any two piles. The player taking the last stone wins.

Find $\displaystyle\sum_{\substack{(a,b,c) \in \mathcal{P} \\ 0 \le a \le b \le c \le 1000}} (a+b+c)$ where $\mathcal{P}$ is the set of P-positions (losing positions for the player to move).

\section{Game-Theoretic Analysis}

\subsection{Sprague--Grundy Framework}

This is an impartial game. Positions are classified as:
\begin{itemize}
    \item \textbf{P-position}: every move leads to an N-position (losing for mover).
    \item \textbf{N-position}: some move leads to a P-position (winning for mover).
\end{itemize}

\subsection{Available Moves from $(a,b,c)$}

\begin{enumerate}
    \item Remove $k \ge 1$ stones from one pile.
    \item Remove $k \ge 1$ stones from two piles simultaneously.
\end{enumerate}

All resulting positions are re-sorted so $a' \le b' \le c'$.

\subsection{Computation}

Process positions in order of increasing $a+b+c$. Position $(0,0,0)$ is a P-position. For each position, check all successors:

\begin{itemize}
    \item If any successor is a P-position, current position is N-position.
    \item If no successor is a P-position, current position is P-position.
\end{itemize}

\section{Optimizations}

\begin{enumerate}
    \item Store P-position status in a 3D boolean array.
    \item Early termination: as soon as we find one move to a P-position, mark as N and move on.
    \item Process in canonical order $a \le b \le c$.
\end{enumerate}

\section{Complexity}

Total positions: $\binom{1001+2}{3} \approx 1.67 \times 10^8$. With early termination and cache-friendly access, computation completes in reasonable time.

\section{Result}

\[
\boxed{167542057}
\]

\section{Answer}

\[
\boxed{167542057}
\]

\end{document}
